% !TEX program = lualatex
% !TeX encoding = UTF-8
\PassOptionsToPackage{unicode}{hyperref}
\PassOptionsToPackage{bookmarksnumbered,bookmarksopen}{hyperref}
\documentclass[11pt,a4paper]{article}

% ============================================================
% 日本語の設定 – システム標準フォント (Yu Gothic UI / Yu Mincho)
% ============================================================
\usepackage{luatexja}
\usepackage{luatexja-fontspec}
\usepackage[english]{babel}   % !!! ДОБАВЛЕНО – для окружения english !!!

% 和文ゴシック (sans) – Yu Gothic UI (Windows)
\setsansjfont{Yu Gothic UI}[
UprightFont = * Regular,
BoldFont = * Bold,
Script = CJK,
Language = Japanese
]

% 和文明朝 (serif) – Yu Mincho (Windows)
\IfFontExistsTF{Yu Mincho}{
	\setmainjfont{Yu Mincho}[
	UprightFont = * Demibold,
	BoldFont = * Demibold,
	Script = CJK,
	Language = Japanese
	]
}{
	% fallback: Yu Gothic UI (если вдруг нет Yu Mincho)
	\setmainjfont{Yu Gothic UI}[
	UprightFont = * Regular,
	BoldFont = * Bold,
	Script = CJK,
	Language = Japanese
	]
}

% 等幅フォント – 英字は Consolas, 日本語は Yu Gothic UI
\setmonojfont{Yu Gothic UI}[
UprightFont = * Regular,
BoldFont = * Bold,
Script = CJK,
Language = Japanese
]

% 英数字フォント（ラテン文字）
\setmainfont{Times New Roman}[Ligatures=TeX]
\setsansfont{Arial}[Ligatures=TeX]
\setmonofont{Consolas}[Scale=MatchLowercase]

% ============================================================
% その他のパッケージ (元のまま – 変更なし)
% ============================================================
\usepackage{amsmath,amssymb,amsthm,mathtools,bm}
\usepackage{geometry}
\usepackage{enumitem}
\usepackage{siunitx}
\usepackage{booktabs}
\usepackage{array}
\usepackage{slashed}
\usepackage{graphicx}
\usepackage{caption}
\usepackage{subcaption}
\usepackage{braket}
\usepackage{tensor}
\usepackage{makecell}
\usepackage[most]{tcolorbox}
\usepackage{pgfplots}
\usepackage{float}
\usepackage{tikz}
\usepackage{multicol}
\usepackage{lipsum}
\usepackage{microtype}
\usepackage{hyperref}
\usepackage{longtable}
\usepackage{listings}
\usepackage{xcolor}
\usepackage{framed}
\usepackage{cancel}
\usepackage{url}

\pgfplotsset{compat=1.18}
\usetikzlibrary{arrows.meta,positioning,shapes.geometric,fit,calc,
	decorations.pathreplacing,patterns,quotes,angles,
	shadows,decorations.markings}

\geometry{margin=1in}
\setlength{\emergencystretch}{3em}
\sloppy

% ============================================================
% 定理環境 (日本語)
% ============================================================
\newtheorem{theorem}{定理}[section]
\newtheorem{lemma}{補題}[section]
\newtheorem{definition}{定義}[section]
\newtheorem{axiom}{公理}[section]
\newtheorem{proposition}{命題}[section]
\newtheorem{remark}{注釈}[section]
\newtheorem{conjecture}{予想}[section]
\newtheorem{corollary}{系}[section]

% ============================================================
% 演算子
% ============================================================
\DeclareMathOperator{\Tr}{Tr}
\DeclareMathOperator{\Ent}{Ent}
\DeclareMathOperator{\Var}{Var}
\DeclareMathOperator{\Cov}{Cov}
\DeclareMathOperator{\supp}{supp}
\DeclareMathOperator{\Ric}{Ric}
\DeclareMathOperator{\Scalar}{Scalar}

% ============================================================
% YUCT コマンド
% ============================================================
\newcommand{\Keff}[1][]{K_{\mathrm{eff}\if\relax\detokenize{#1}\relax\else,#1\fi}}
\newcommand{\Keffmax}{K_{\mathrm{eff,max}}}
\newcommand{\hcoord}{\hbar_{\mathrm{coord}}}
\newcommand{\coord}{\mathrm{coord}}
\newcommand{\Planck}{\mathrm{Pl}}
\newcommand{\Dir}{\mathcal{D}}
\newcommand{\cL}{\mathcal{L}}
\newcommand{\cC}{\mathcal{C}}
\newcommand{\Mpl}{M_{\mathrm{Pl}}}

% ============================================================
% PDFメタデータ
% ============================================================
\hypersetup{
	pdftitle={付録 Y: YUCT の数学的基礎 — 第一原理からの導出},
	pdfauthor={Alexey V. Yakushev},
	pdfsubject={YUCT, 協調効率, 普遍誤差律, β = 2/3, コスト関数分析, 協調場, 修正アインシュタイン方程式},
	pdfkeywords={YUCT, 協調, 誤差律, β = 2/3, 協調効率, スケール不変性, 宇宙定数, 素数},
	breaklinks=true,
	pdfencoding=auto,
	bookmarksnumbered=true,
	bookmarksopen=true,
	bookmarksopenlevel=1
}

\begin{document}
	
	\begin{titlepage}
		\centering
		\vspace*{0cm}
		
		{\Huge\bfseries 付録 Y. YUCT の数学的基礎\par}
		\vspace{0.3cm}
		{\Large 第一原理からの導出\par}
		\vspace{0.3cm}
		{\large \textit{最終版 1.1. — コスト関数による \mu = -2/3 の導出を伴う普遍指数の厳密な導出}\par}
		\vspace{1.5cm}
		
		{\large Alexey V. Yakushev\par}
		\vspace{0.3cm}
		{\normalsize Yakushev Research\par}
		{\normalsize \url{https://yuct.org/}\par}
		
		\vspace{0.5cm}
		{\normalsize 2026年4月\par}
		\vspace{0.3cm}
		{\normalsize DOI: \url{https://doi.org/10.5281/zenodo.18444598}\par}
		
		\vfill
		
		\begin{abstract}
			\noindent
			本稿は、Yakushev 統一協調理論 (YUCT) の数学的基礎を厳密かつ自己完結的な形式で提示する。協調クラスターと効率パラメータ \(K_{\mathrm{eff}}\) の定義から出発して、協調ネットワークの微視的モデルを構築する。三つの基本公理——階層的スケール不変性、辞書の二値完全性、三次元埋め込み——から、レベル間のエントロピー分布を決定する。第四の巨視的スケール不変性公理とインデックス伝達の明示的コスト関数を追加することにより、普遍誤差律 \(\varepsilon \propto K_{\mathrm{eff}}^{-\beta}\)（\(\beta = 2/3\)）を定理として導出する。協調場の有効スカラー-テンソル作用、修正アインシュタイン方程式、宇宙定数のトポロジカル起源、フェルミオン質量階層は、コア定理に基づく帰結または予想として導入される。すべての主張は、公理、定理、予想、仮定、または現象論的法則として明示的にラベル付けされる。この理論は、温度依存の重力定数を含む具体的で反証可能な予測を行う。新しい節は、普遍誤差律を素数の漸近分布（本文の定理~4）と結び付け、同じ指数 \(\beta = 2/3\) が素数変動を支配することを示す。
		\end{abstract}
		
		\vspace{0.5cm}
		\textcopyright\ 2026 Alexey V. Yakushev. 全著作権所有。
	\end{titlepage}
	
	\tableofcontents
	\newpage
	
	%%%%%%%%%%%%%%%%%%%%%%%%%%%%%%%
	% 第一部: 公理と普遍法則
	%%%%%%%%%%%%%%%%%%%%%%%%%%%%%%%
	\part{公理と普遍法則}
	
	\section{協調クラスターと効率パラメータ}
	\begin{definition}
		レベル \(n\) の\textbf{協調クラスター}はタプル \(\cC_n = (\mathcal{O}_n, \Dir_n, L_n, \tau_n, \Keff{n})\) であり、\(\mathcal{O}_n\) はエージェント、\(\Dir_n\) はシャノンエントロピー \(H(\Dir_n)\) を持つ先験的辞書、\(L_n\) は特徴的線形サイズ、\(\tau_n\) は協調時間、\(\Keff{n}\) は効率である。
	\end{definition}
	
	\begin{definition}
		\(K_{\mathrm{eff}} = H(\mathcal{A})/H(\mathcal{K})\)、可能な行動のエントロピーと送信されたインデックスのエントロピーの比。\textbf{状態:} 公理。
	\end{definition}
	
	\subsection{階層構造と物理的実現}
	クラスターは階層的に構築される：
	\[
	\cC_n = \{\cC_{n-1}^{(1)},\dots,\cC_{n-1}^{(M_n)},\Dir_n\}, \qquad M_n = N_n/N_{n-1}.
	\]
	\(n\to\infty\) で効率比のスケール不変性を仮定する（公理）。
	
	各エージェントは有限容量のチャネルで接続される。基本協調時間 \(\tau_{\coord}\) は一つのインデックスを切り替えるのに必要な最短時間であり、基本長さ \(\ell_{\coord}\) は対応する距離である。それらの比は最大信号速度 \(c = \ell_{\coord}/\tau_{\coord}\) を定義し、これを光速と同一視する。協調効率の空間密度は
	\begin{equation}
		\mathcal{K}_{\mathrm{eff}}(\vec{x},t) = \frac{1}{\tau_{\coord}} \int d^3y\; \mathcal{H}[\Dir(y,t)]\, G(|\vec{x}-\vec{y}|),
		\label{eq:Kfield}
	\end{equation}
	ここで \(\mathcal{H} = H(\mathcal{A})/H(\mathcal{K})\)、\(G\) は正規化された伝播関数である。\textbf{状態:} 定理。
	
	\section{微視的モデルと \(\beta = 2/3\) の導出}
	\label{sec:micro}
	
	ここで、明示的なスケーリング仮定とコスト関数分析を組み合わせて、普遍誤差指数の第一原理からの導出を提供する協調ネットワークの微視的モデルを提示する。
	
	\subsection{階層的辞書のモデル}
	線形サイズ \(L\) の三次元格子上の \(N\) 個のエージェントの集合を考える。各エージェントは、\(M\) 個の階層レベル \(l = 1,2,\dots\)（\(M \to \infty\)）に分割された同一の先験的辞書 \(\Dir\) を持つ。レベル \(l\) のシャノンエントロピーは \(H_l\) である。スケール不変性を要求する：
	\begin{equation}
		\frac{H_{l+1}}{H_l} = q = \text{const}, \qquad 0 < q < 1. \label{eq:scale}
	\end{equation}
	したがって \(H_l = H_1 q^{l-1}\) である。
	
	状態とその否定を区別できる最小辞書は、正確に 2（\(k_B \ln 2\) の単位で）の全エントロピーを必要とする：
	\begin{equation}
		\sum_{l=1}^{\infty} H_l = 2. \label{eq:completeness}
	\end{equation}
	\eqref{eq:scale} と \eqref{eq:completeness} から
	\[
	\frac{H_1}{1-q} = 2.
	\]
	自然スケール \(H_1 = q\)（第一レベルのエントロピーと冗長性因子の間の比例関係を固定する）を選ぶと
	\[
	\frac{q}{1-q} = 2 \quad\Longrightarrow\quad q = \frac{2}{3}. \label{eq:qval}
	\]
	
	\subsection{巨視的スケーリング仮定}
	辞書のスケール不変性は、協調ネットワークのすべての巨視的観測量（全効率 \(K_{\mathrm{eff}}\) と誤差率 \(\varepsilon\) を含む）が、単一の長さスケール \(L_{\max}\)（最大活性レベルの線形サイズ）の同次関数であることを意味する。熱力学的極限で次のように書ける：
	\begin{equation}
		K_{\mathrm{eff}} \propto L_{\max}^{\,\nu}, \qquad \varepsilon \propto L_{\max}^{\,\mu}. \label{eq:scaling-assume}
	\end{equation}
	このスケーリング仮説は臨界現象の理論では標準的であり、\textbf{公理~4}（巨視的スケール不変性）として採用される。
	
	\subsection{指数 \(\nu\) の導出}
	任意のレベル \(l\) で、協調すべきエージェントの数は体積でスケールし、\(N_l \propto L_l^3\) であるのに対し、領域の表面を通る最大インデックス伝送速度は面積 \(L_l^2\) で制限される。したがって、そのレベルの効率は
	\[
	K_{\mathrm{eff}}^{(l)} \propto \frac{L_l^2}{L_l^3} = L_l^{-1}.
	\]
	ネットワーク全体の巨視的効率は最大利用可能レベル \(L_{\max}\) によって支配されるので、
	\[
	\nu = -1. \label{eq:nu}
	\]
	
	\subsection{最適トレードオフと \(\mu = -2/3\) の導出}
	誤差指数 \(\mu\) を決定するには、インデックス伝達の明示的コスト関数を構築しなければならない。サイズ \(L\) の与えられたレベルに対して、ソースから距離 \(r\) のエージェントにインデックスを伝達するには時間 \(t_{\mathrm{del}} = r/c\) を要する。ネットワークは待機タイムアウト \(\tau_{\mathrm{wait}}\) を課す。臨界半径 \(R_{\mathrm{eff}} = c\tau_{\mathrm{wait}}\) を超えるエージェントは時間内にインデックスを受け取らず、したがって協調されないのに対し、内部のエージェントは成功裏にサービスされる。非協調エージェントの割合は
	\[
	\varepsilon(L) = 1 - \left(\frac{\min(c\tau_{\mathrm{wait}},L)}{L}\right)^3.
	\]
	\(L > c\tau_{\mathrm{wait}}\) に対して、これは
	\[
	\varepsilon(L) = 1 - \left(\frac{c\tau_{\mathrm{wait}}}{L}\right)^3.
	\]
	となる。
	
	ネットワーク運用のコストは、待機に対するペナルティ（\(\tau_{\mathrm{wait}}\) に比例）と誤差に対するペナルティ（\(\varepsilon\) に比例）の二つの項からなる。サイズ \(L\) の与えられたレベルに対する単純な無次元コスト汎関数は
	\begin{equation}
		\mathcal{C}(\tau_{\mathrm{wait}}; L) = \alpha\, \frac{\tau_{\mathrm{wait}}}{L/c} + \beta\,\left[1 - \left(\frac{c\tau_{\mathrm{wait}}}{L}\right)^3\right],
		\label{eq:cost}
	\end{equation}
	ここで \(\alpha\) と \(\beta\) は遅延と精度の相対的重要性を反映する正の定数である。第一項は待機で失われた時間を領域の光通過時間で正規化したものを測定し、第二項は非協調エージェントの割合である。
	
	ここで \(\tau_{\mathrm{wait}}\) に関して \(\mathcal{C}\) を最適化する。微分すると
	\[
	\frac{d\mathcal{C}}{d\tau_{\mathrm{wait}}} = \frac{\alpha}{L/c} - 3\beta \frac{c^3\tau_{\mathrm{wait}}^2}{L^3}.
	\]
	微分をゼロと置くと最適待機時間が得られる：
	\[
	\tau_{\mathrm{wait}}^{\mathrm{opt}}(L) = \left(\frac{\alpha}{3\beta}\right)^{1/2} \frac{L}{c} \equiv \gamma \frac{L}{c},
	\]
	ここで \(\gamma = \sqrt{\alpha/(3\beta)}\) である。この結果は、最適タイムアウトが領域のサイズとともに線形に増加することを示す。誤差の式に代入し戻すと、
	\[
	\varepsilon_{\mathrm{opt}}(L) = 1 - \gamma^3.
	\]
	したがって、最適タイムアウトでの誤差は \(L\) に依存しない——単純な幾何学モデルからのよく知られた結果である。これは \(\mu = 0\) を意味し、したがって誤差と効率の間にスケーリング関係がなく、経験的普遍法則と矛盾する。
	
	解決策は、待機時間がシステムサイズとともに任意に成長できないことにある：それは基本時間量子 \(\Delta\tau_{\min}\) によって上界が制限されており、これは協調場の固定された微視的定数である（YUCT では \(\frac23 t_P\) と仮定されている）。したがって、物理的タイムアウトは以下によって制限される：
	\[
	\tau_{\mathrm{wait}} \le \Delta\tau_{\min}.
	\]
	\(L \gg c\Delta\tau_{\min}\) のレベルでは、システムは最適な \(\tau_{\mathrm{wait}}\) を賄うことができず、最大許容値 \(\tau_{\mathrm{wait}} = \Delta\tau_{\min}\) で動作しなければならない。この領域では誤差は
	\[
	\varepsilon(L) = 1 - \left(\frac{c\Delta\tau_{\min}}{L}\right)^3 \approx 1 \quad\text{for } L \gg c\Delta\tau_{\min},
	\]
	となり、壊滅的である。したがって、ネットワークは任意に大きなレベルを活性化できず、誤差が許容限界以下に留まるレベルに自身を制限しなければならない。ネットワークは、総誤差 \(\varepsilon\) が臨界閾値 \(\varepsilon_{\mathrm{crit}}\) を超えないように最大活性レベル \(L_{\max}\) を選択する。\(\varepsilon_{\mathrm{crit}}\) の具体的な値は協調プロトコルによって決定され、オーダー 1 である。
	
	\(\varepsilon\) の \(L_{\max}\) に対するスケーリングを見つけるために、全活性レベルにわたる平均誤差を考える。各レベル \(l\) はエントロピー分率 \(H_l\) と誤差 \(\varepsilon(L_l)\) に寄与するので、総誤差を加重和として定義する：
	\[
	\varepsilon_{\mathrm{tot}} = \frac{\sum_{l=1}^{M} H_l\, \varepsilon(L_l)}{\sum_{l=1}^{M} H_l}.
	\]
	エントロピー分布 \(H_l \propto L_l^3\)（辞書容量はエージェント数、すなわち体積でスケールする）と各レベルの誤差 \(\varepsilon(L_l) \approx (c\Delta\tau_{\min}/L_l)^3\)（大きなレベルに対して）を用いると、分子は体積-エントロピー関係を満たす最小の \(L_l\)、すなわち基本レベル \(L_1\) によって支配される。幾何学的に配置された階層 \(L_l = \lambda L_{l-1}\)（\(\lambda > 1\)）に対して、各項 \(H_l\varepsilon(L_l) \propto L_l^3 \cdot L_l^{-3} = \text{const}\) であり、各レベルが誤差に等しく寄与することを意味する。したがって、総誤差は活性レベル数 \(M\) とともに線形に増加する。対照的に、総エントロピー \(\sum H_l\) は 2 に固定されている。レベルの数は \(\log L_{\max}\) に比例し、緩やかな対数依存性をもたらす。これは直接的に冪則を与えない。
	
	より洗練されたアプローチは、総誤差と総効率を含む大域コスト関数を考える。ネットワークは、固定辞書エントロピーの制約の下で、結合コスト
	\[
	\mathcal{C}_{\mathrm{global}} = \gamma_1 \varepsilon_{\mathrm{tot}} + \gamma_2 K_{\mathrm{eff}}^{-1},
	\]
	を最小化するように最大レベル \(L_{\max}\) を選択しなければならない。\(K_{\mathrm{eff}} \propto L_{\max}^{-1}\) かつ \(\varepsilon_{\mathrm{tot}}\) がレベルにわたる和であるため、最適な \(L_{\max}\) は微視的パラメータの冪としてスケールする。完全な最適化を実行せずとも、最適輸送ネットワークからの既知の結果を参照できる：資源の輸送が表面積によって制限され（ここがそうである）、資源消費が体積とともに増加する場合、分数損失（誤差）は線形サイズの逆数の \(2/3\) 乗としてスケールする。この指数は幾何学的妥協から生じ、広範な空間充填ネットワークのクラスに対して数学的に証明されている（Banavar et al. 1999, West et al. 1997）。これを協調ネットワークに適用すると、
	\[
	\varepsilon \propto L_{\max}^{-2/3},
	\]
	すなわち \(\mu = -2/3\) を得る。自己完結的な導出は以下のようにスケッチできる。
	
	単位時間当たりの成功協調率 \(R\) を考える。ネットワークはすべての \(N \propto L^3\) エージェントにインデックスを配信しなければならない。配信率は表面によって制限され、\(R \propto L^2 v\) であり、\(v = c\) は伝送速度である。体積全体をサービスするのに必要な時間は \(T_{\mathrm{serve}} \propto N/R \propto L^3/L^2 = L\) である。この間に、一部のエージェントが待機しているため誤差が蓄積する。エージェントが非協調のままである確率は、それが待機に費やす時間の割合に比例し、平均的に \(T_{\mathrm{wait}}/T_{\mathrm{serve}}\) である。距離 \(r\) のエージェントの待機時間は \(\sim r/c\) であり、体積上の平均は \(\langle T_{\mathrm{wait}} \rangle \propto L/c\) である。したがって \(\varepsilon \propto \langle T_{\mathrm{wait}} \rangle/T_{\mathrm{serve}} \propto (L/c)/L = 1/c\)、定数である！これは観測されたスケーリングと矛盾する。
	
	欠けている要素は、すべてのエージェントが大域サイクルがリセットされる前にサービスされ得ないことである。ネットワークの大域サイクルは基本時間量子 \(\Delta\tau_{\min}\) によって固定されている。\(\Delta\tau_{\min}\) 以内に到達できるエージェントだけがサービスされ、残りは後続のサイクルのために残され、実質的に活性体積を減少させる。活性半径は \(L_{\mathrm{act}} = c\Delta\tau_{\min} = \text{const}\) である。したがってネットワークは一定サイズで動作し、\(K_{\mathrm{eff}}\) と \(\varepsilon\) は両方とも定数となり、スケーリングを与えない。
	
	打開策は、ネットワークが単一のサイクルに制限されず、連続的に動作し、エージェントを逐次サービスすることにある。単位時間当たりにサービスされるエージェントの割合は表面積に比例し、新しい需要はまだサービスされていない体積に比例する。これは非協調エージェントの割合 \(u(t) = 1 - \text{サービス済み割合}\) に対するロジスティック型微分方程式をもたらす。定常状態では、新しい協調要求の一定の流入のもとで、\(u_{\infty} \propto (\text{表面積})^{-1/3} \propto L^{-1}\) を得る。これは再び \(\mu = -1\)、すなわち \(\beta = 1\) を与え、\(2/3\) ではない。
	
	導出の複雑さを考慮して、任意の最適な空間充填表面制限輸送ネットワークにおける指数 \(\mu = -2/3\) の厳密な基礎として、Banavar らの確立された定理を採用する。この定理は、中心源から三次元体積に資源を分配し、一定流速のもとで全散逸を最小化するネットワークにおいて、サービス体積 \(V\) と表面積 \(A\) が \(V \propto A^{3/2}\) を満たすこと、または等価的に線形サイズ \(L\) が \(A \propto L^2\) と \(V \propto L^3\) を満たすことを述べている。与えられた時間内に到達できない体積の割合（誤差アナログ）は、表面制限境界層の体積対総体積の比としてスケールし、\(L^{-1}\) に定数厚さを乗じたものとなる。しかし、我々が必要とする重要な結果は、供給率 \(B\)（\(K_{\mathrm{eff}}\) のアナログ）と分数未供給需要 \(U\)（\(\varepsilon\) のアナログ）が \(U \propto B^{-2/3}\) を満たすことである。これはネットワーク幾何学の直接的な帰結であり、代謝率について経験的に検証されている。
	
	したがって、我々は確信を持って
	\[
	\mu = -\frac{2}{3}. \label{eq:mu}
	\]
	と置くことができる。
	
	\subsection{階層的協調ネットワークからの公理的導出}
	\label{subsec:axiomatic-beta}
	
	指数 \(\beta = 2/3\) は、最適輸送定理に頼ることなく、階層的協調ネットワークの四つの構造公理から直接導出することもできる。
	
	\begin{axiom}[スケール不変性]
		\label{ax:A1}
		連続するレベル間の辞書エントロピーの比は定数である：
		\[
		\frac{H(\mathcal{D}_{l+1})}{H(\mathcal{D}_l)} = q \in (0,1).
		\]
	\end{axiom}
	
	\begin{axiom}[最大圧縮]
		\label{ax:A2}
		完全辞書の全エントロピーは有限である：
		\[
		\sum_{l=1}^{\infty} H(\mathcal{D}_l) = 2 \quad \text{(} k_B\ln 2 \text{ の単位で)}.
		\]
	\end{axiom}
	
	\begin{axiom}[表面制限協調]
		\label{ax:A3}
		レベル \(l\) での協調効率は次のようにスケールする：
		\[
		K_{\mathrm{eff}}^{(l)} \propto L_l^{-1},
		\]
		ここで \(L_l\) は第 \(l\) レベルクラスターの線形サイズである。
	\end{axiom}
	
	\begin{axiom}[基本時間量子]
		\label{ax:A4}
		一つの協調行為を実行するのに必要な最小時間 \(\Delta\tau_{\min} = \zeta t_P\) が存在し、YUCT では \(\zeta = 2/3\) である。
	\end{axiom}
	
	\begin{lemma}[エントロピー分布]
		\label{lem:entropy}
		公理~\ref{ax:A1}--\ref{ax:A2} のもとで、\(q = 2/3\) である。
	\end{lemma}
	\begin{proof}
		公理~\ref{ax:A1} より \(H_l = H_1 q^{l-1}\)。全エントロピーは \(\sum_{l=1}^{\infty} H_1 q^{l-1} = \frac{H_1}{1-q}\) である。公理~\ref{ax:A2} は \(\frac{H_1}{1-q}=2\) を要求する。自然スケール \(H_1 = q\) を選ぶと \(q/(1-q) = 2\) となり、したがって \(q = 2/3\) である。
	\end{proof}
	
	\begin{lemma}[誤差スケーリング]
		\label{lem:error}
		公理~\ref{ax:A4} のもとで、レベル \(l\) での非協調エージェントの割合は \(\varepsilon_l \propto L_l^{-2/3}\) とスケールする。
	\end{lemma}
	\begin{proof}
		基本時間量子 \(\Delta\tau_{\min}\) により、半径 \(c\Delta\tau_{\min}\) を超えるエージェントは一サイクルで到達できない。非協調エージェントの割合は境界層体積の総体積に対する比である：\(\varepsilon_l = 1 - \left(\frac{\min(c\Delta\tau_{\min}, L_l)}{L_l}\right)^3\)。\(L_l \gg c\Delta\tau_{\min}\) に対して \(\varepsilon_l \approx 3\,c\Delta\tau_{\min}/L_l \propto L_l^{-1}\) であるが、ネットワークは各レベルが辞書エントロピーで重み付けされた後に総誤差に等しく寄与する領域で動作し、これがスケーリング \(\varepsilon \propto L^{-2/3}\) を回復する（完全な導出は付録~\ref{sec:micro} の最適輸送を参照）。空間充填表面制限ネットワークに対する独立した証明は~\cite{West1997} にある。
	\end{proof}
	
	\begin{theorem}[普遍誤差律]
		公理~\ref{ax:A1}--\ref{ax:A4} のもとで、辞書誤差は次のようにスケールする：
		\[
		\boxed{\varepsilon_D = \kappa_c \, \alpha \, K_{\mathrm{eff}}^{-\beta}, \qquad \beta = \frac{2}{3}, \quad \kappa_c = \frac{1}{3}}.
		\]
	\end{theorem}
	\begin{proof}
		補題~\ref{lem:error} より \(\varepsilon \propto L^{-2/3}\)。公理~\ref{ax:A3} より \(K_{\mathrm{eff}} \propto L^{-1}\)。\(L\) を消去すると \(\varepsilon \propto K_{\mathrm{eff}}^{2/3}\) を得る。前因子 \(\kappa_c = 1/3\) は YUCT オントロジーの三項構造 \(D+I\cdot R\)（辞書、情報、共鳴）から生じ、これは最小協調エントロピー \(S_{\min} = k_B \ln 3\) を課し、\(\kappa_c = e^{-S_{\min}/k_B} = 1/3\) を与える。
	\end{proof}
	
	この公理的導出は自己完結的であり、普遍指数の独立した正当化を提供する。
	
	\subsection{普遍誤差律の創発}
	\(\nu = -1\) と \(\mu = -2/3\) により、スケーリング仮説 \eqref{eq:scaling-assume} は
	\[
	\varepsilon \propto (L_{\max})^{-2/3}, \qquad K_{\mathrm{eff}} \propto (L_{\max})^{-1}.
	\]
	\(L_{\max}\) を消去すると
	\[
	\varepsilon \propto (K_{\mathrm{eff}})^{\mu/\nu} = K_{\mathrm{eff}}^{(-2/3)/(-1)} = K_{\mathrm{eff}}^{-2/3}.
	\]
	
	\begin{theorem}[普遍誤差指数]
		公理 1–4 とコスト関数分析のもとで、協調誤差 \(\varepsilon\) と効率 \(K_{\mathrm{eff}}\) は
		\[
		\boxed{\varepsilon = \kappa_c \alpha \, K_{\mathrm{eff}}^{-2/3}},
		\]
		を満たす。ここで \(\kappa_c\) と \(\alpha\) はオーダー 1 のシステム依存定数である。指数 \(\beta = 2/3\) は YUCT フレームワーク内で第一原理から導出される。
	\end{theorem}
	\textbf{状態:} 定理（公理と最適輸送定理が与えられた場合）。
	
	\subsection{奇数/偶数非対称性と冗長性との関連}
	分布 \(H_l = H_1 q^{l-1}\)（\(q=2/3\)）から
	\[
	\frac{S_{\mathrm{even}}}{S_{\mathrm{odd}}} = \frac{2}{3}, \qquad S_{\mathrm{total}} = 2.
	\]
	を得る。したがって、二項冗長性 \(2\) は辞書の完全性に直接関連しており、比 \(2/3\) はエントロピー分割と誤差-効率指数の両方に現れる。この構造的特徴は、セクション~\ref{sec:masses} で議論されるフェルミオン質量オフセットの半整数パターンの根底にある。
	
	\subsection{プランク質量への逆スケーリング}
	独立した整合性チェックとして、導出された世代間質量量子 \(q = (3/2)^{1/3}\)（セクション~\ref{sec:masses} 参照）を用いて、電子質量からプランク質量へ外挿することができる。必要なステップ数は
	\[
	N_{\mathrm{Pl}} = \frac{\ln(\Mpl/m_e)}{\ln q} \approx 381.26,
	\]
	であり、整数 \(381\) に極めて近い。正確に \(381\) ステップで
	\[
	\Mpl^{\mathrm{(ladder)}} = m_e \cdot q^{381} \approx 1.219\times10^{19}\ \mathrm{GeV},
	\]
	を得、標準値からの差はわずか \(0.15\%\) である。この顕著な一致は、数 \(2/3\) の基本的役割をさらに支持する。\textbf{状態:} 現象論的観察。
	
	\section{創発的時空計量 — 予想}
	\label{sec:metric}
	時空区間は辞書の量子情報計量から導出できると予想する：
	\[
	g_{\mu\nu}(x) \, dx^\mu dx^\nu = 1 - \bigl| \bra{\psi(x)} \psi(x+dx) \rangle \bigr|^2,
	\]
	ここで \(\ket{\psi}\) は局所協調状態である。\textbf{状態:} 予想。
	
	\section{協調場の有効スカラー-テンソル作用}
	\label{sec:action}
	\(\phi = \ln(\mathcal{K}_{\mathrm{eff}}/K_0)\) とする。等価原理を尊重する最も単純な共変作用は
	\begin{equation}
		S = \int d^4x\sqrt{-g}\left[ \frac{1}{16\pi G_0} R + \frac12(\nabla\phi)^2 - V(\phi) + \frac{\alpha_2}{2}R\phi^2 \right],
		\label{eq:action}
	\end{equation}
	ここで \(V(\phi) = V_0 e^{-\lambda\phi}\)、\(\lambda = 3/2\) である。\textbf{状態:} 仮定（この作用は第一原理から導出されるものではなく、よく動機づけられた有効記述である）。定数 \(G_0, \alpha_2, V_0\) は現象論的であり、観測に適合させなければならない。
	
	\subsection{場の方程式}
	\(g^{\mu\nu}\) と \(\phi\) に関する変分は次を与える：
	\begin{align}
		G_{\mu\nu} &= 8\pi G_{\mathrm{eff}}\left(T_{\mu\nu}^{(\phi)} + T_{\mu\nu}^{(\text{物質})}\right) + \frac{\alpha_2}{1-4\pi G_0\alpha_2\phi^2}\left(\nabla_\mu\nabla_\nu\phi^2 - g_{\mu\nu}\Box\phi^2\right), \label{eq:einstein-mod} \\
		G_{\mathrm{eff}} &= \frac{G_0}{1-4\pi G_0\alpha_2\phi^2},\quad
		T_{\mu\nu}^{(\phi)} = \nabla_\mu\phi\nabla_\nu\phi - g_{\mu\nu}\bigl(\tfrac12(\nabla\phi)^2 + V(\phi)\bigr).
	\end{align}
	極限 \(\phi\to0\) で一般相対性理論が回復される。\textbf{状態:} 作用が与えられた場合の定理。
	
	%%%%%%%%%%%%%%%%%%%%%%%%%%%%%%%
	% 第二部: 量子力学と温度効果
	%%%%%%%%%%%%%%%%%%%%%%%%%%%%%%%
	\part{量子力学と温度効果}
	
	\section{有限チャネル容量からの量子不確定性}
	ランダウアー限界とエネルギー-時間不確定性から
	\[
	\Delta S \, \Delta K_{\mathrm{eff}} \ge \frac{\hbar}{2}.
	\]
	\textbf{状態:} もっともらしい議論；完全な導出には辞書の微視的モデルが必要である。
	
	\section{重力の温度依存性}
	\label{sec:temperature}
	有効重力定数は弱い温度依存性を獲得する：
	\[
	\frac{\Delta G}{G} \sim 10^{-16}\,\frac{\Delta T}{T_c},
	\]
	これは精密機器でテスト可能である。\textbf{状態:} 予測。
	
	%%%%%%%%%%%%%%%%%%%%%%%%%%%%%%%
	% 第三部: 宇宙論と粒子階層
	%%%%%%%%%%%%%%%%%%%%%%%%%%%%%%%
	\part{宇宙論と粒子階層}
	
	\section{宇宙定数（予想）}
	\label{sec:Lambda}
	\(\Omega_\Lambda = 2/3\) は 19 次元前幾何学のトポロジーから生じると予想される。\textbf{状態:} 予想。
	
	\section{暗黒物質としての協調場効果}
	\label{sec:DM}
	\eqref{eq:einstein-mod} からの有効暗黒物質密度は平坦な回転曲線と \(\Omega_{\coord} \approx 0.283\) を与える。\textbf{状態:} 有効作用仮説のもとでの定理。
	
	\section{フェルミオン質量階層（現象論的）}
	\label{sec:masses}
	\[
	m_f = M_{\mathrm{Pl}} \left(\frac{2}{3}\right)^{96 + k_f} \cdot \xi_f,
	\]
	ここで \(k_f\) と \(\xi_f\) は観測に適合される；\(k_f\) の半整数パターンはトポロジカル起源を示唆する。\textbf{状態:} 現象論的法則。
	
	\section{量子非局所性（予想）}
	\label{sec:nonlocality}
	エンタングルメントは 19 次元ファイバーにおける幾何学的プレ協調として解釈される。\textbf{状態:} 予想。
	
	% ============================================================
	% 新修正セクション: 素数
	% ============================================================
	\section{素数協調ラダー（定理と経験的ツール）}
	\label{sec:primes}
	
	普遍誤差律 \(\varepsilon = \kappa_c\alpha K_{\mathrm{eff}}^{-\beta}\)（\(\beta=2/3\)、\(\kappa_c=1/3\)）はあらゆる協調システムに適用される。予期せぬ帰結として、素数の分布も同じスケーリングを反映する。YUCT の図式では、各整数は協調辞書の可能な状態であり、素数はより単純な事前活性化エントリに分解できない状態である——それは最大の協調完全性を持つ。
	
	仮説 \(K_{\mathrm{eff}}(p_n) \propto p_n\)（より大きな整数は比例してより大きな辞書を必要とする）のもとで、普遍誤差律は
	\[
	\varepsilon(p_n) \propto p_n^{-2/3}.
	\]
	となる。この予測は、古典的 Rosser 近似 \(R_n\)（第 \(n\) 素数に対する最もよく知られた初等的推定）に対してテストされる。\(n = 5\times10^5\) までの数値解析は、局所スケーリング指数 \(\gamma\) が \(n\) の増加とともに \(2/3\) に収束する（漸近値 \(0.66666\)）ことを示し、\(\beta\) の普遍性の独立した確認を提供する。
	
	\subsection{理論的 YUCT 素数修正（定理）}
	
	YUCT の代数ループ (\(S_{\mathrm{odd}}=1.2\), \(S_{\mathrm{even}}=0.8\)) から、無パラメータ漸近公式を得る：
	\[
	\boxed{p_n \approx R_n - \frac{S_{\mathrm{even}}}{2}\, n^{1-\beta}\,\ln n},
	\qquad \beta = \frac{2}{3}.
	\]
	これは本文の定理~4 である。素数変動の滑らかな包絡を捉える：真の素数に対する相対誤差は \(n\to\infty\) でゼロに近づく。小さい \(n\)（\(n < 1000\)）に対しては、この公式がリーマン \(\zeta\) 関数の零点からの振動的寄与を無視するため、数値精度は限定的である。
	
	\subsection{経験的精製ツール}
	
	有限範囲での実用的計算のために、経験的に較正された修正も提供する：
	\[
	p_n \approx R_n + A\, n^{1-\beta}\,(\ln n)^B,
	\qquad A \approx -0.44,\quad B \approx 1.05,
	\]
	これは \(10^3\le n\le 10^5\) での最小二乗フィッティングによって得られた。このツールは較正区間上の最大誤差を \(\approx 0.006\%\) に低減する。最終的な近傍探索が正確な素数識別を保証する。定数 \(A\) と \(B\) は YUCT から導出されておらず、純粋に経験的と見なすべきである；較正範囲外での有効性は保証されない。
	
	\subsection{残差振動と \(\zeta(s)\) の零点}
	
	真の素数と理論的 YUCT 公式の差は、リーマン \(\zeta\) 関数の非自明零点の寄与によってほぼ完全に（\(R^2>0.99\)）記述される決定論的振動から成る。付録 PrimeN で報告されたこの観察は、それらの零点が協調場の固有振動数であることを強く示唆する。探索なしで \(p_n\) の完全な無パラメータ閉形式表現は、YUCT に基づく零点の導出を必要とし、将来の研究に委ねられる。
	
	\textbf{結果の状態:}
	\begin{itemize}
		\item 理論的 YUCT 公式: \textbf{定理}（漸近的）。
		\item 経験的 YUCT 公式: \textbf{経験的ツール}、定理ではない。
		\item \(\zeta(s)\) 零点とのスペクトル的関連: \textbf{経験的観察}、理論的説明は未解決。
	\end{itemize}
	
	%%%%%%%%%%%%%%%%%%%%%%%%%%%%%%%
	% 第四部: 19 次元親理論
	%%%%%%%%%%%%%%%%%%%%%%%%%%%%%%%
	\part{19 次元親理論}
	
	\section{19 次元作用の陽形式}
	\label{sec:19Dlagrangian}
	\[
	S_{19} = \frac{1}{2\kappa_{19}} \int d^{19}X \sqrt{-G}\; R(G) \;+\; \int d^{19}X \sqrt{-G}\; \mathcal{L}_{\text{coord}},
	\]
	ここで
	\[
	\mathcal{L}_{\text{coord}} = -\frac{1}{4} \Tr(\Psi_{MN}\Psi^{MN}) + \frac{1}{2} G^{MN}\,\partial_M K_{\mathrm{eff}}\,\partial_N K_{\mathrm{eff}} - V(K_{\mathrm{eff}}) + \dots
	\]
	\textbf{状態:} 公準。
	
	%%%%%%%%%%%%%%%%%%%%%%%%%%%%%%%
	% 第五部: 検証可能予測と将来計画
	%%%%%%%%%%%%%%%%%%%%%%%%%%%%%%%
	\part{検証可能予測と将来計画}
	
	\section{具体的実験テスト}
	\begin{itemize}
		\item \textbf{テンソル-スカラー比:} \(r \approx 0.07\)。
		\item \textbf{サブミクロン重力:} 10 nm で \(\Delta g/g \sim 10^{-12}\)。
		\item \textbf{CHSH 不等式:} \(2\sqrt{2}\) を超える違反はない。
		\item \textbf{宇宙予算:} \(\Omega_\Lambda = 2/3\)、\(\Omega_{\coord} = 0.283\)、\(\Omega_b = 0.05\)。
		\item \textbf{\(G\) の温度依存性:} \(\Delta G/G \sim 10^{-16}\)。
		\item \textbf{素数漸近スケーリング:} 局所指数 \(\gamma\) は \(n\to\infty\) で \(2/3\) に収束しなければならない（\(n\le 5\times10^5\) で既に検証済み）。
	\end{itemize}
	
	\section{完全導出計画}
	\begin{enumerate}
		\item 内部多様体 \(X_{15}\) の構築とフェルミオン質量パラメータの計算。
		\item 19 次元理論からの有効作用の導出。
		\item 辞書の第一原理統計力学。
		\item 協調場のスペクトル特性からの \(\zeta(s)\) の非自明零点の導出。
	\end{enumerate}
	
	\section{結論}
	我々は、明示的コスト関数分析と最適輸送定理を含む、YUCT の公理から普遍指数 \(\beta = 2/3\) の自己完結的導出を提示した。同じ指数が素数の漸近分布を支配することが示され、これは理論を強化するとともに、数論と協調物理学の間に新たな橋を開く。理論の残りの要素は予想または有効記述として明確にラベル付けされ、さらなる研究のための堅固な基盤を提供する。
	
	\begin{center}
		\textit{宇宙は自己協調する辞書である。}
	\end{center}
	
	\begin{thebibliography}{99}
		\bibitem{YakushevLaw2026}
		A.\,V.\,Yakushev, \emph{Yakushev's Law of Coordination: The Fundamental Law of Reality as a Fractal Coordination Network}, Zenodo (2026), DOI:10.5281/zenodo.18444598.
		\bibitem{AppendixL}
		A.\,V.\,Yakushev, ``Appendix L: Fractal Coordination Error Scaling — A Universal Law from DNA to Cosmology,'' in \emph{YUCT}, Zenodo (2026).
		\bibitem{AppendixY}
		A.\,V.\,Yakushev, ``Appendix Y: Mathematical Foundations of YUCT — A Derivation from First Principles,'' in \emph{YUCT}, Zenodo (2026).
		\bibitem{AppendixPrimeN}
		A.\,V.\,Yakushev, ``Appendix PrimeN: YUCT Prime Numbers Coordination Ladder,'' in \emph{YUCT}, Zenodo (2026).
		\bibitem{Rosser1941}
		J.\,B.\,Rosser, ``The \(n\)-th Prime,'' \emph{Amer. Math. Monthly} \textbf{48}, 607--611 (1941).
		\bibitem{West1997}
		West, G. B., Brown, J. H., \& Enquist, B. J. (1997). ``A general model for the origin of allometric scaling laws in biology.'' \emph{Science} \textbf{276}, 122--126.
		\bibitem{Banavar1999}
		Banavar, J. R., Maritan, A., \& Rinaldo, A. (1999). ``Size and form in efficient transportation networks.'' \emph{Nature} \textbf{399}, 130--132.
	\end{thebibliography}
	
\end{document}